\documentclass[10pt,a4paper]{amsart}
\usepackage{amsmath,amssymb,amsthm,graphicx}
% Retained preamble for the Stacks Project chapter snapshot in this directory.
% The chapter's own preamble.tex is a copy of the shared Stacks Project preamble; the block
% below is sliced MECHANICALLY, line for line, from it. Nothing was added, removed or
% redefined; the wrapper class, the xr-hyper/hyperref wiring and the externaldocument list
% are omitted because the assembler supplies the document class.
%
% lines taken from preamble.tex: 3, 5, 9, 11, 12, 156, 157, 158, ... (51 lines total)
%
% package coverage restored from the chapter's own preamble (2026-10-13): xy, xy-2cell

\usepackage[all]{xy}
\xyoption{2cell}
\UseAllTwocells

\newenvironment{reference}{\comment}{\endcomment}
\newenvironment{slogan}{\comment}{\endcomment}
\newenvironment{history}{\comment}{\endcomment}
\theoremstyle{plain}
\newtheorem{theorem}[subsection]{Theorem}
\newtheorem{proposition}[subsection]{Proposition}
\newtheorem{lemma}[subsection]{Lemma}
\theoremstyle{definition}
\newtheorem{definition}[subsection]{Definition}
\newtheorem{example}[subsection]{Example}
\newtheorem{exercise}[subsection]{Exercise}
\newtheorem{situation}[subsection]{Situation}
\theoremstyle{remark}
\newtheorem{remark}[subsection]{Remark}
\newtheorem{remarks}[subsection]{Remarks}
\numberwithin{equation}{subsection}
\def\lim{\mathop{\mathrm{lim}}\nolimits}
\def\colim{\mathop{\mathrm{colim}}\nolimits}
\def\Spec{\mathop{\mathrm{Spec}}}
\def\Hom{\mathop{\mathrm{Hom}}\nolimits}
\def\Ext{\mathop{\mathrm{Ext}}\nolimits}
\def\SheafHom{\mathop{\mathcal{H}\!\mathit{om}}\nolimits}
\def\SheafExt{\mathop{\mathcal{E}\!\mathit{xt}}\nolimits}
\def\Sch{\mathit{Sch}}
\def\Mor{\mathop{\mathrm{Mor}}\nolimits}
\def\Ob{\mathop{\mathrm{Ob}}\nolimits}
\def\Sh{\mathop{\mathit{Sh}}\nolimits}
\def\NL{\mathop{N\!L}\nolimits}
\def\CH{\mathop{\mathrm{CH}}\nolimits}
\def\proetale{{pro\text{-}\acute{e}tale}}
\def\etale{{\acute{e}tale}}
\def\QCoh{\mathit{QCoh}}
\def\Ker{\mathop{\mathrm{Ker}}}
\def\Im{\mathop{\mathrm{Im}}}
\def\Coker{\mathop{\mathrm{Coker}}}
\def\Coim{\mathop{\mathrm{Coim}}}
\DeclareMathSymbol{\boxtimes}{\mathbin}{AMSa}{"02}
\def\QCohstack{\mathcal{QC}\!\mathit{oh}}
\def\Cohstack{\mathcal{C}\!\mathit{oh}}
\def\Spacesstack{\mathcal{S}\!\mathit{paces}}
\def\Quotfunctor{\mathrm{Quot}}
\def\Hilbfunctor{\mathrm{Hilb}}
\def\Curvesstack{\mathcal{C}\!\mathit{urves}}
\def\Polarizedstack{\mathcal{P}\!\mathit{olarized}}
\def\Complexesstack{\mathcal{C}\!\mathit{omplexes}}
\def\Pic{\mathop{\mathrm{Pic}}\nolimits}
\def\Picardstack{\mathcal{P}\!\mathit{ic}}
\def\Picardfunctor{\mathrm{Pic}}
\def\Deformationcategory{\mathcal{D}\!\mathit{ef}}
\usepackage{cleveref}
\crefname{originalequation}{original equation}{original equations}
\Crefname{originalequation}{Original equation}{Original equations}
\makeatletter
\newlabel{equation-triangle}{{08QS}{source}{}{sourcepointer-1}{}}
\newlabel{equation-triangle@cref}{{[originalequation][0][]08QS}{[1][0][]source}{}{}{}}
\makeatother
\crefname{originallemma}{original lemma}{original lemmas}
\Crefname{originallemma}{Original lemma}{Original lemmas}
\makeatletter
\newlabel{lemma-cotangent-complex-polynomial-algebra}{{08QH}{source}{}{sourcepointer-2}{}}
\newlabel{lemma-cotangent-complex-polynomial-algebra@cref}{{[originallemma][0][]08QH}{[1][0][]source}{}{}{}}
\makeatother
\crefname{originallemma}{original lemma}{original lemmas}
\Crefname{originallemma}{Original lemma}{Original lemmas}
\makeatletter
\newlabel{lemma-mod-Koszul-regular-ideal}{{08SK}{source}{}{sourcepointer-3}{}}
\newlabel{lemma-mod-Koszul-regular-ideal@cref}{{[originallemma][0][]08SK}{[1][0][]source}{}{}{}}
\makeatother
\crefname{originallemma}{original lemma}{original lemmas}
\Crefname{originallemma}{Original lemma}{Original lemmas}
\makeatletter
\newlabel{lemma-relation-with-naive-cotangent-complex}{{08RB}{source}{}{sourcepointer-4}{}}
\newlabel{lemma-relation-with-naive-cotangent-complex@cref}{{[originallemma][0][]08RB}{[1][0][]source}{}{}{}}
\makeatother
\crefname{originallemma}{original lemma}{original lemmas}
\Crefname{originallemma}{Original lemma}{Original lemmas}
\makeatletter
\newlabel{more-algebra-lemma-quasi-regular-ideal-finite-projective}{{08RK}{source}{}{sourcepointer-5}{}}
\newlabel{more-algebra-lemma-quasi-regular-ideal-finite-projective@cref}{{[originallemma][0][]08RK}{[1][0][]source}{}{}{}}
\makeatother
\title{The cotangent complex of a local complete intersection map is perfect with tor amplitude in [-1, 0]}
\author{The Stacks Project Authors}
\date{stacks-project:cotangent.tex:local-snapshot-20261012; GNU FDL 1.2 or later}
\begin{document}
\maketitle
\setcounter{subsection}{1}
\noindent\textit{Editorial note.} All delivered mathematical text is the author's own, copied from the pinned Stacks Project chapter file cotangent.tex; no statement or proof step was rewritten, reordered or removed.\par\medskip
\noindent\textit{Editorial note.} The statement is two sentences; the proof separates the perfectness statement from the concentration of cohomology and closes with the distinguished triangle.\par\medskip
\noindent\textit{Editorial note.} The five references of the counted spans are recorded as non-load-bearing navigation pointers; the argument that consumes them stays inside the counted proof span.\par\medskip
\noindent\textit{Editorial note.} Preamble: the retained preamble of this chapter (sources/stacks-cotangent/retained-preamble.tex, lines 1-63), containing the Xy-pic setup, the comment environments, the theorem declarations and the macros including \textbackslash{}NL.\par\medskip
\noindent\textit{Allowed standard prerequisites (editorial).}\begin{itemize}
\item The identification of the naive cotangent complex with the truncated cotangent complex. Lemma lemma-relation-with-naive-cotangent-complex, used at proof lines 2965-2968. Turns the two-term presentation complex into the truncation.
\item Finite projectivity of the conormal module of a quasi-regular ideal. More on Algebra, Lemma more-algebra-lemma-quasi-regular-ideal-finite-projective, used at proof lines 2969-2973. Makes the truncation a perfect complex with tor amplitude in [-1, 0].
\item The distinguished triangle for the composition A -> P -> B. equation-triangle, used at proof lines 2975-2981. Reduces the vanishing of the remaining cohomology groups to that of L\_\{B/P\}.
\item The Koszul regularity of the ideal in the local complete intersection case and the polynomial algebra computation. Lemmas lemma-mod-Koszul-regular-ideal and lemma-cotangent-complex-polynomial-algebra, used at proof lines 2980-2983. Give the concentration in degrees -1 and 0 and the computation of the left term.
\end{itemize}\medskip
\ifdefined\hypertarget\hypertarget{sourcepointer-1}{}\fi
\noindent\textit{Source reference.} Original equation 08QS: focused-seeds/AGP/sources/stacks-snapshot-20261011/tags.txt, source lines 11070--11070. Navigation pointer into the Stacks Project; the standard fact it supplies is applied inside the counted proof span and the argument resting on it stays there.\par\medskip
\ifdefined\hypertarget\hypertarget{sourcepointer-2}{}\fi
\noindent\textit{Source reference.} Original lemma 08QH: focused-seeds/AGP/sources/stacks-snapshot-20261011/tags.txt, source lines 11060--11060. Navigation pointer into the Stacks Project; the standard fact it supplies is applied inside the counted proof span and the argument resting on it stays there.\par\medskip
\ifdefined\hypertarget\hypertarget{sourcepointer-3}{}\fi
\noindent\textit{Source reference.} Original lemma 08SK: focused-seeds/AGP/sources/stacks-snapshot-20261011/tags.txt, source lines 11133--11133. Navigation pointer into the Stacks Project; the standard fact it supplies is applied inside the counted proof span and the argument resting on it stays there.\par\medskip
\ifdefined\hypertarget\hypertarget{sourcepointer-4}{}\fi
\noindent\textit{Source reference.} Original lemma 08RB: focused-seeds/AGP/sources/stacks-snapshot-20261011/tags.txt, source lines 11089--11089. Navigation pointer into the Stacks Project; the standard fact it supplies is applied inside the counted proof span and the argument resting on it stays there.\par\medskip
\ifdefined\hypertarget\hypertarget{sourcepointer-5}{}\fi
\noindent\textit{Source reference.} Original lemma 08RK: focused-seeds/AGP/sources/stacks-snapshot-20261011/tags.txt, source lines 11098--11098. Navigation pointer into the Stacks Project; the standard fact it supplies is applied inside the counted proof span and the argument resting on it stays there.\par\medskip
\begin{proposition}
\label{proposition-cotangent-complex-local-complete-intersection}
Let $A \to B$ be a local complete intersection map.
Then $L_{B/A}$ is a perfect complex with tor amplitude in $[-1, 0]$.
\end{proposition}

\begin{proof}
Choose a surjection $P = A[x_1, \ldots, x_n] \to B$ with kernel $J$.
By Lemma \ref{lemma-relation-with-naive-cotangent-complex}
we see that $J/J^2 \to \bigoplus B\text{d}x_i$
is quasi-isomorphic to $\tau_{\geq -1}L_{B/A}$.
Note that $J/J^2$ is finite projective
(More on Algebra, Lemma
\ref{more-algebra-lemma-quasi-regular-ideal-finite-projective}),
hence $\tau_{\geq -1}L_{B/A}$ is a perfect complex with
tor amplitude in $[-1, 0]$.
Thus it suffices to show that $H^i(L_{B/A}) = 0$ for $i \not \in [-1, 0]$.
This follows from (\ref{equation-triangle})
$$
L_{P/A} \otimes_P^\mathbf{L} B \to L_{B/A} \to L_{B/P} \to
L_{P/A} \otimes_P^\mathbf{L} B[1]
$$
and Lemma \ref{lemma-mod-Koszul-regular-ideal}
to see that $H^i(L_{B/P})$ is zero unless $i \in \{-1, 0\}$.
(We also use Lemma \ref{lemma-cotangent-complex-polynomial-algebra}
for the term on the left.)
\end{proof}

\end{document}
